
\subsubsection{Automation}\label{trien-khai-remediate-automation}

Khắc phục tự động được thực thi qua các playbook Ansible để cấu hình nhanh chóng và nhất quán các quy chuẩn bảo mật trên toàn bộ hệ thống. Các playbook remediation tập trung vào việc chỉnh sửa cấu hình Docker daemon, phân quyền file, cấu hình container runtime, network, resource limit và dọn dẹp tài nguyên dư thừa.

\textbf{Bảng chi tiết các hoạt động Remediation tự động:}

{\def\LTcaptype{none} % do not increment counter
\begin{longtable}{@{}|
  >{\raggedright\arraybackslash}p{(\linewidth - 8\tabcolsep - 5\arrayrulewidth) * \real{0.0973}}|
  >{\raggedright\arraybackslash}p{(\linewidth - 8\tabcolsep - 5\arrayrulewidth) * \real{0.1936}}|
  >{\raggedright\arraybackslash}p{(\linewidth - 8\tabcolsep - 5\arrayrulewidth) * \real{0.3545}}|
  >{\raggedright\arraybackslash}p{(\linewidth - 8\tabcolsep - 5\arrayrulewidth) * \real{0.3545}}|@{}}
\hline
\begin{minipage}[b]{\linewidth}\raggedright
\textbf{Nhóm section}
\end{minipage} & \begin{minipage}[b]{\linewidth}\raggedright
\textbf{Nội dung Remediation Auto}
\end{minipage} & \begin{minipage}[b]{\linewidth}\raggedright
\textbf{Phương thức tự động hóa / Module sử dụng}
\end{minipage} & \begin{minipage}[b]{\linewidth}\raggedright
\textbf{Kết quả cần đạt sau remediation}
\end{minipage} \\
\hline
\endhead

\multirow{2}{=}{1} & 1.1.3 - 1.1.18: Cấu hình Audit rules tự động. & Module \texttt{ansible.builtin.apt}, \texttt{ansible.builtin.copy}, \texttt{ansible.builtin.lineinfile} và \texttt{ansible.builtin.service}. Playbook cài đặt \texttt{auditd}, ghi các rule theo dõi Docker daemon, Docker socket, containerd và các file cấu hình liên quan. & Các audit rules được ghi vào hệ thống; dịch vụ \texttt{auditd} được bật và khởi động thành công; các thay đổi liên quan đến Docker được ghi log để phục vụ kiểm toán. \\
\cline{2-4}

& 1.2.2: Tự động cập nhật Docker. & Module \texttt{ansible.builtin.apt}; cập nhật package index và nâng cấp Docker Engine nếu có phiên bản mới trong repository. & Docker Engine được cập nhật lên phiên bản phù hợp, giúp giảm rủi ro từ các lỗ hổng bảo mật đã được vá. \\
\hline

\multirow{6}{=}{2} & 2.2 - 2.3: Hạn chế traffic giữa container và đặt logging level phù hợp. & Module \texttt{ansible.builtin.copy}, \texttt{ansible.builtin.template} hoặc \texttt{ansible.builtin.lineinfile}; chỉnh file \texttt{/etc/docker/daemon.json} để cấu hình \texttt{icc} và \texttt{log-level}. & Inter-container communication trên default bridge được hạn chế nếu chính sách yêu cầu; logging level được đặt về \texttt{info}, tránh dùng \texttt{debug} trong production. \\
\cline{2-4}

& 2.5: Loại bỏ insecure registry không an toàn. & Module \texttt{ansible.builtin.lineinfile} hoặc \texttt{ansible.builtin.template}; xóa hoặc không khai báo trường \texttt{insecure-registries} trong \texttt{daemon.json}. & Docker daemon không sử dụng registry không mã hóa hoặc không đáng tin cậy; image chỉ được pull từ nguồn được phê duyệt. \\
\cline{2-4}

& 2.6 - 2.7: Sử dụng storage driver phù hợp. & Module \texttt{ansible.builtin.template}; cấu hình \texttt{storage-driver} trong \texttt{daemon.json}. Với môi trường lab có thể đặt \texttt{overlay2}; trong production cần backup dữ liệu trước khi thay đổi. & Docker không dùng \texttt{aufs} hoặc \texttt{devicemapper} không phù hợp; storage driver được cấu hình theo khuyến nghị, ưu tiên \texttt{overlay2}. \\
\cline{2-4}

& 2.9 - 2.10: Cấu hình default ulimit và user namespace. & Module \texttt{ansible.builtin.template}; thêm \texttt{default-ulimits} và \texttt{userns-remap} vào \texttt{daemon.json}. & Container có giới hạn tài nguyên mặc định hợp lý; user namespace remapping được bật nếu chính sách hệ thống yêu cầu. \\
\cline{2-4}

& 2.14 - 2.15: Bật no-new-privileges và live restore. & Module \texttt{ansible.builtin.template}; cấu hình \texttt{no-new-privileges} và \texttt{live-restore} trong \texttt{daemon.json}, sau đó restart Docker bằng handler. & Container không được leo thang đặc quyền mới; container có thể tiếp tục chạy khi Docker daemon restart. \\
\cline{2-4}

& 2.17 - 2.19: Tắt userland proxy, giữ seccomp và tắt experimental features. & Module \texttt{ansible.builtin.template} hoặc \texttt{ansible.builtin.lineinfile}; cấu hình \texttt{userland-proxy}, kiểm soát seccomp profile và \texttt{experimental}. & \texttt{userland-proxy} được tắt nếu không cần thiết; seccomp không bị disable; experimental features không bật trong môi trường production. \\
\hline

\multirow{6}{=}{3} & 3.1 - 3.2: Quyền file service Docker. & Module \texttt{ansible.\allowbreak builtin.\allowbreak file}; xác định service file bằng \texttt{systemctl show} rồi đặt owner/group và mode phù hợp. & File service thuộc sở hữu root; permission là \texttt{0644} hoặc hạn chế hơn. \\
\cline{2-4}

& 3.3 - 3.4: Khắc phục ownership và permission của \texttt{docker.socket}. & Module \texttt{ansible.builtin.file}; xác định đường dẫn \texttt{docker.socket} và đặt lại owner/group/mode. & File \texttt{docker.socket} thuộc sở hữu \texttt{root:root}; permission là \texttt{0644} hoặc hạn chế hơn. \\
\cline{2-4}

& 3.5 - 3.6: Khắc phục quyền thư mục \texttt{/etc/docker}. & Module \texttt{ansible.builtin.file}; đặt owner/group và mode cho thư mục cấu hình Docker. & Thư mục \texttt{/etc/docker} thuộc sở hữu \texttt{root:root}; permission là \texttt{0755} hoặc hạn chế hơn. \\
\cline{2-4}

& 3.7 - 3.10: Khắc phục quyền file chứng chỉ registry và Docker TLS. & Module \texttt{ansible.builtin.find} kết hợp \texttt{ansible.builtin.file}; duyệt các file trong \texttt{/etc/docker/certs.d/} và các file certificate/key nếu tồn tại. & Certificate thuộc sở hữu \texttt{root:root}; private key có permission chặt chẽ, không để user không liên quan đọc hoặc ghi. \\
\cline{2-4}

& 3.15 - 3.16: Khắc phục quyền Docker socket thực tế. & Module \texttt{ansible.builtin.file}; đặt owner/group/mode cho \texttt{/var/run/docker.sock}. & Docker socket thuộc sở hữu \texttt{root:docker}; permission là \texttt{0660}, hạn chế truy cập trái phép vào Docker daemon. \\
\cline{2-4}

& 3.17 - 3.18: Khắc phục quyền \texttt{daemon.json}, sysconfig Docker và containerd socket. & Module \texttt{ansible.builtin.file}; kiểm tra file tồn tại rồi đặt owner/group/mode phù hợp. & File cấu hình Docker và containerd socket có ownership đúng; permission không quá rộng, giảm nguy cơ bị chỉnh sửa trái phép. \\
\hline

\multirow{6}{=}{4} & 4.1: Cấu hình container chạy bằng non-root user. & Chỉnh Dockerfile hoặc cấu hình container bằng Ansible; thêm chỉ thị \texttt{USER} hoặc cấu hình user khi tạo container. & Container không chạy bằng user \texttt{root}, giảm rủi ro leo thang đặc quyền khi container bị khai thác. \\
\cline{2-4}

& 4.5: Bật Docker Content Trust. & Cấu hình biến môi trường \texttt{DOCKER\_CONTENT\_TRUST=1} trong môi trường build/deploy hoặc pipeline CI/CD. & Docker chỉ pull/run image có chữ ký hợp lệ nếu hệ thống áp dụng Content Trust. \\
\cline{2-4}

& 4.6: Thêm chỉ thị \texttt{HEALTHCHECK}. & Chỉnh Dockerfile hoặc cấu hình \texttt{healthcheck} khi tạo container bằng module Docker của Ansible. & Image hoặc container có healthcheck để Docker theo dõi trạng thái ứng dụng. \\
\cline{2-4}

& 4.7: Không để lệnh update đứng độc lập trong Dockerfile. & Chỉnh Dockerfile, gộp lệnh update và install trong cùng một lớp build, đồng thời dọn cache package sau khi cài đặt. & Dockerfile không có lệnh update độc lập gây cache sai hoặc bỏ sót bản vá. \\
\cline{2-4}

& 4.9: Dùng \texttt{COPY} thay cho \texttt{ADD}. & Chỉnh Dockerfile; thay các lệnh \texttt{ADD} không cần thiết bằng \texttt{COPY}. & Dockerfile dùng \texttt{COPY} cho thao tác sao chép file thông thường, tránh hành vi ngầm của \texttt{ADD}. \\
\cline{2-4}

& 4.10 - 4.12: Không đưa secrets vào Dockerfile và xác minh package/artifact. & Kết hợp Ansible, CI/CD và secret manager; đưa secrets ra biến môi trường hoặc vault, kiểm tra chữ ký package/artifact trước khi build. & Dockerfile không chứa password, token, private key; package và artifact có nguồn gốc đáng tin cậy. \\
\hline

\multirow{9}{=}{5} & 5.9, 5.14: Giới hạn port và interface. & \begin{minipage}[t]{\linewidth}\raggedright
\texttt{community.docker.}\\
\texttt{docker\_container} với \texttt{published\_ports}.\strut
\end{minipage} & Port chỉ bind vào interface được phê duyệt, không mở rộng không cần thiết ra toàn bộ host. \\
\cline{2-4}

& 5.10, 5.30: Tạo và sử dụng network riêng. & \begin{minipage}[t]{\linewidth}\raggedright
\texttt{community.docker.}\\
\texttt{docker\_network} và\\
\texttt{community.docker.}\\
\texttt{docker\_container}.\strut
\end{minipage} & Container dùng \texttt{secure-net}, không dùng host network hoặc default bridge. \\
\cline{2-4}

& 5.11, 5.12, 5.19, 5.29: Giới hạn tài nguyên. & Thuộc tính \texttt{memory}, \texttt{cpu\_shares}, \texttt{ulimits}, \texttt{pids\_limit}. & Container có giới hạn RAM, CPU priority, file descriptor và PID. \\
\cline{2-4}

& 5.13: Root filesystem read-only. & \texttt{read\_only: true} và các mount \texttt{tmpfs}. & Root filesystem không ghi được; ứng dụng chỉ ghi vào vùng được phép. \\
\cline{2-4}

& 5.15: Giới hạn restart. & \texttt{restart\_policy} và \texttt{restart\_retries}. & \texttt{on-failure} với tối đa 5 lần retry. \\
\cline{2-4}

& 5.26: Chặn leo thang đặc quyền. & \texttt{security\_opts}. & Có \texttt{no-new-privileges}. \\
\cline{2-4}

& 5.27: Thiết lập healthcheck. & Thuộc tính \texttt{healthcheck}. & Healthcheck tồn tại và container đạt trạng thái \texttt{healthy}. \\
\cline{2-4}

& 5.28: Ghim image. & Biến \texttt{image\_name} dùng version tag hoặc digest. & Không dùng \texttt{latest} hoặc image không tag. \\
\cline{2-4}

& 5.4: Thu hẹp capability. & \texttt{cap\_drop} và chỉ add capability cần thiết. & Không giữ capability mặc định không cần thiết. \\
\hline

6 & 6.1 - 6.2: Tự động dọn dẹp tài nguyên dư thừa. & Module \texttt{ansible.builtin.cron} hoặc chạy prune tự động bằng lệnh Docker CLI. & Giải phóng thành công image, container, volume hoặc network không còn sử dụng; đĩa sạch và không tích tụ rác. \\
\hline

\end{longtable}
}
```
